\subsection{時間変化格子と ALE 効果（本稿では未実装）}
\label{app:grid_ale}
% ═══════════════════════════════════════════════════════════════════════════════

本付録は §\ref{sec:grid_gen} の格子が時間発展する場合の ALE（Arbitrary Lagrangian-Eulerian）補正を説明する．
\textbf{本稿の実装では格子を時間固定として扱う（$\bm{v}_\text{mesh} = 0$）}か，タイムステップごとに格子を再生成するが $\bm{v}_\text{mesh}$ の補正は省略する．

ALE 定式化では格子速度 $\bm{v}_\text{mesh}$ を輸送方程式に追加する必要がある：$\partial\psi/\partial t + \nabla\cdot((\bm{u} - \bm{v}_\text{mesh})\psi) = 0$．ただしこの追加項の離散化は実装の複雑さを大きく増加させる．

\textbf{誤差の発生機構：}
1ステップあたりの界面移動量を $\delta = |\bm{u}|\Delta t$ とすると，
ALE 補正を省略することで $\bm{v}_\text{mesh}$ による追加対流量が無視される：
\[
  \varepsilon_{\text{ALE}} \sim |\bm{v}_\text{mesh}| \cdot |\bnabla\psi| \cdot \Delta t \sim U_\Gamma \cdot \frac{1}{\varepsilon} \cdot \Delta t
\]
ここで $|\bnabla\psi| \sim 1/(4\varepsilon)$（界面近傍プロファイル），$U_\Gamma$ は界面速度である．

\textbf{近似が有効な条件：}
\begin{center}
\renewcommand{\arraystretch}{1.3}
\begin{tabular}{>{\raggedright\arraybackslash}p{48mm}c>{\raggedright\arraybackslash}p{42mm}}
\hline
条件 & 移動量 $\delta/h$ & ALE 省略の妥当性 \\
\hline
緩やかな界面変形 & $< 0.1$ & 良好（誤差 $< 10\%$ of $\Delta t^2$ 項） \\
標準的な問題 & $\approx 0.5$ & 条件付きで可（格子収束確認要） \\
大密度比 ($\rho_l/\rho_g{=}1000$) の高速界面 & $> 1$ & 精度劣化の可能性大（ALE 実装推奨） \\
\hline
\end{tabular}
\end{center}

$\rho_l/\rho_g = 1000$ の場合，液相変形が速く $\delta/h > 1$ となることがあり，ALE 補正なしの格子更新（Mode 2）は実質的に $O(\Delta t)$ 精度（1次）となりうる．
本稿の実装では ALE 補正を省略しており，上昇気泡ベンチマーク（§\ref{sec:bench_bubble}）では時間刻み精度の確認を格子収束テストとあわせて実施することを推奨する．

\textbf{Mode 1（時間固定格子，本稿の主設計）：}
格子速度 $\bm{v}_\text{mesh} = 0$ として格子を更新しない．ALE 誤差はゼロであり，全体時間精度 $O(\Delta t^2)$ が保持される（§\ref{sec:time_accuracy_table}）．
ただし，格子は初期水準集合値 $\phi^{(0)}$ から生成されるため，界面が初期精細化領域から大きく離れる問題（例：上昇気泡ベンチマーク，§\ref{sec:bench_bubble}）では経時的に空間解像度が失われる．
長時間・長距離移動を伴う計算では，適切な間隔で格子を再生成する（Mode 2 + ALE 実装）か，一様格子を選択することを推奨する．

\textbf{Mode 2（タイムステップごとに格子を再生成，ALE 省略）：}
上表の $\delta/h > 1$ 条件に該当する場合，$O(\Delta t)$ に精度が劣化し，本稿の全体時間精度 $O(\Delta t^2)$ の主張（§\ref{sec:time_accuracy_table}）と矛盾する．
$\rho_l/\rho_g = 1000$（上昇気泡系）では特に注意が必要であり，Mode 2 を採用する場合は ALE 補正（$\bm{v}_\text{mesh} \neq 0$）の実装を強く推奨する（本稿では未実装）．

% ═══════════════════════════════════════════════════════════════════════════════
